\usepackage[nomaketitle]{config/psl-cover/psl-cover}

\sujet{\huge Universalité dans des modèles de populations branchants, structurés et régulés}
\entitle{\Large Universality in structured and regulated branching population models}
\author{Mathilde André}

\encadrant{Directeur \textsc{De Recherche}}

\institute{l'École Normale Supérieure}
\doctoralschool{École Doctorale de Mathématiques de Paris Centre}{386}
\specialty{Mathématiques}
\laboratory{IBENS, École Normale Supérieure\\ CIRB, Collège de France}
\date{15 décembre 2026}

\pslassetspath{config/psl-cover}
% \jurymember{1}{Prénom \textsc{Nom 1}}{Affiliation & \emph{Examinateur}}{} % le président doit être en premier
\jurymember{2}{Vincent \textsc{Bansaye}}{Professeur, École Polytechnique}{Rapporteur}
\jurymember{3}{Sarah \textsc{Penington}}{Reader and Royal Society Fellow,\\University of Bath}{Rapportrice}
\jurymember{4}{Raphaël \textsc{Forien}}{CR, Inrae PACA}{Examinateur}
\jurymember{5}{Hélène \textsc{Leman}}{CR Inria, École Normale Supérieure de Lyon}{Examinatrice}
\jurymember{6}{Amaury \textsc{Lambert}}{Professeur, École Normale Supérieure et Collège de France}{Directeur}
\jurymember{7}{Emmanuel \textsc{Schertzer}}{Professor, University of Vienna}{Directeur}
\jurymember{8}{Jean-Jil \textsc{Duchamps}}{MCF, Université Marie et Louis Pasteur}{Co-encadrant}
% \jurymember{7}{Prénom NOM}{Titre, établissement}{Directeur de thèse}

% entre 1700 et 4000 signes ? 
\frabstract{
À mesure que les populations évoluent, se reproduisent, font face à des catastrophes et s'adaptent à de nouveaux environnements, des mutations s'accumulent dans leur ADN. L'un des enjeux centraux de la génétique des populations consiste à reconstituer certains aspects de cette histoire complexe à partir des échantillons génétiques, souvent peu nombreux et incomplets, dont on dispose. Cette thèse propose une approche probabiliste pour comprendre la diversité génétique au sein de modèles de populations structurées et en interaction. Le fil rouge est de trouver des limites d'échelles dans lesquelles des objets universels et des comportements plus simples émergent de de dynamiques écologiques complexes.

La thèse comporte trois chapitres, ainsi qu'une section consacrée à des problèmes ouverts. Les deux premiers chapitres sont consacrés à l'étude des lignées ancestrales dans des modèles de populations régulés. Le Chapitre~\ref{chp3_singletype} développe une approche fondée sur les moments pour analyser la généalogie de processus de branchement densité-dépendants. Nous utilisons ce cadre pour montrer que, au voisinage d'un équilibre écologique stable, la généalogie d'un échantillon converge vers le coalescent de Kingman, lorsque la taille typique de la population tend vers l'infini. Le Chapitre~\ref{chp4_multitype} permet d'étendre la classe d'universalité du coalescent de Kingman aux populations structurées et régulées. Nous développons une méthodologie générale reliant la dynamique de fractions neutres aux descriptions généalogiques en temps rétrospectif. Nous montrons que, malgré la complexité de la structure sous-jacente et des interactions écologiques, les échelles de temps écologique et généalogique se découplent asymptotiquement. Après un changement d'échelle approprié, la généalogie converge à nouveau vers le coalescent neutre standard. Dans le  Chapitre~\ref{chp2_ks_countable}, nous nous intéressons à un autre résultat d'universalité et établissons des résultats de type Kesten--Stigum pour des processus de Galton--Watson dans un espace de types infini dénombrable. Nous identifions une condition nécessaire et suffisante d'uniforme intégrabilité de la martingale additive et de convergence de la distribution des types en probabilité. Nous montrons de nouveaux critère pour la convergence presque sûre, puis appliquons ces résultats pour obtenir des approximations de branchement dans des modèles structurés d'épidémiologie. Le dernier Chapitre~\ref{chp5_opening} étend, à l'aide d'arguments heuristiques, le cadre développé dans les Chapitres~\ref{chp3_singletype} et~\ref{chp4_multitype} afin d'étudier l'émergence de $\Lambda$-coalescents dans des populations structurées et en interaction. Enfin nous y discutons de l'application des résultats du Chapitre~\ref{chp2_ks_countable} à l'étude des distances typiques dans un réseau orienté construit à partir d'un nombre dénombrable de modèles de configuration.
}

\enabstract{
While populations evolve, reproduce, face catastrophes, and adapt to new environments, mutations accumulate in their DNA. A central challenge in population genetics is to reconstruct aspects of a population's complex history from the few and often incomplete genetic samples available. This thesis develops a probabilistic approach to understanding genetic diversity in structured and interacting population models. The common thread is the emergence of universal objects and simpler behaviors from complex ecological dynamics under suitable scaling limits. 

The thesis consists of three chapters and an open problem section. The first two chapters are dedicated to the study of ancestral lineages in regulated populations. Chapter~\ref{chp3_singletype} develops a moment-based framework for analyzing the genealogy of density-dependent branching processes. We use this framework to show that, near a stable ecological equilibrium, the genealogy of a sample converges to the Kingman coalescent in the large-population limit. Chapter~\ref{chp4_multitype} extends the class of universality of the Kingman coalescent to structured populations with interactions. We develop a general methodology that relates the forward dynamics of neutral fractions to the backward-in-time descriptions of genealogies. We show that, despite complex ecological structure and interactions, ecological and genealogical timescales asymptotically decouple. After an appropriate rescaling, the genealogy again converges to the standard neutral coalescent. In Chapter~\ref{chp2_ks_countable}, we study another universality result and establish Kesten--Stigum-type results for Galton--Watson processes with countably many types. Under positive recurrence, we identify a sharp condition for uniform integrability and establish convergence of the type distribution in probability and, under broad conditions, almost surely. We apply these results to derive branching approximations in structured models of epidemics. The last Chapter~\ref{chp5_opening} extends the framework developed in Chapters~\ref{chp3_singletype} and~\ref{chp4_multitype}, through heuristic arguments, towards the emergence of $\Lambda$-coalescents from structured interacting populations. Finally, we discuss how the results of Chapter~\ref{chp2_ks_countable} can be applied to studying typical distances in a directed network constructed from countably many configuration models.
}

\frkeywords{Processus de branchement, Populations structurées, Théorèmes limites, Lignées ancestrales, Densité dépendance, Changement de mesure}
\enkeywords{Branching processes, Structured populations, Limit theorems, Ancestral lineages, Density dependence, Change of measure}


%%%%%%% Paramètres PDF
\title{\sujet}
\hypersetup{
pdftitle={Thesis Year},
pdfsubject={\sujet},
pdfauthor={\theauthor},
pdfkeywords={\thefrkeywords},
}
